% ============================================================
%  profile/cv.tex  --  CV source for linkedin-outreach skill
% ============================================================
%
%  ANNOTATION SYSTEM
%  -----------------
%  Add  % @tag  at the end of any \resumeItem line to mark
%  which audience it targets. The linkedin-outreach agent reads
%  these to select / deselect bullets when adapting per company.
%
%  Current tags:
%    @data_engineer   ETL, Spark, Airflow, dbt, pipelines, warehouses
%    @genai           LLM apps, RAG, fine-tuning, AI in production
%    @llm             LLM-specific work (subset of @genai)
%    @rag             Retrieval-augmented generation (subset of @genai)
%    @ml_engineer     Feature engineering, PyTorch, MLflow, model ops
%    @backend         APIs, Go, Docker, Kubernetes, cloud infra
%    @cost            Cost reduction, perf optimisation, automation gains
%    @leadership      Tech lead, mentoring, cross-team ownership
%
%  No tag = always included (education, key differentiators, languages).
%
%  ONE CV RULE
%  -----------
%  Keep this as the single master file. The agent creates adapted
%  copies at /tmp/cv_adapted_*.tex per target -- you never edit
%  multiple source files manually.
%
%  BUILD
%  -----
%    python scripts/build_cv.py \
%      --input  profile/cv.tex \
%      --output profile/cv_builds/cv_base.pdf
%
%  EXTRA PACKAGES (beyond standard texlive)
%    texlive-fonts-extra  (fontawesome5, noto-sans)
%    MiKTeX: auto-installs on first compile
%    Ubuntu: sudo apt install texlive-fonts-extra
%
% ============================================================

\documentclass[letterpaper,11pt]{article}

\usepackage{latexsym}
\usepackage[empty]{fullpage}
\usepackage{titlesec}
\usepackage{marvosym}
\usepackage[usenames,dvipsnames]{color}
\usepackage{enumitem}
\usepackage[colorlinks=true,urlcolor=primary,linkcolor=primary]{hyperref}
\usepackage{fancyhdr}
\usepackage{tabularx}
\usepackage{fontawesome5}
\usepackage[sfdefault]{noto-sans}

\definecolor{primary}{RGB}{0,79,144}

\pagestyle{fancy}
\fancyhf{}
\renewcommand{\headrulewidth}{0pt}

\addtolength{\oddsidemargin}{-0.5in}
\addtolength{\textwidth}{0.92in}
\addtolength{\topmargin}{-0.72in}
\addtolength{\textheight}{1.45in}

\setlength{\tabcolsep}{0in}
\raggedbottom
\raggedright

\titleformat{\section}{
  \vspace{-2pt}\scshape\raggedright\large\color{primary}
}{}{0em}{}[\color{primary}\titlerule \vspace{-4pt}]

\newcommand{\resumeItem}[1]{\item\small{#1}}

\newcommand{\resumeSubheading}[4]{
  \vspace{2pt}\item
  \begin{tabular*}{0.97\textwidth}{l@{\extracolsep{\fill}}r}
    \textbf{#1} & #2 \\
    \textit{\small #3} & \textit{\small #4} \\
  \end{tabular*}\vspace{-5pt}
}

\newcommand{\resumeProjectHeading}[2]{
  \item\vspace{-1pt}
  \begin{tabular*}{0.97\textwidth}{l@{\extracolsep{\fill}}r}
    \small #1 & #2 \\
  \end{tabular*}\vspace{-2pt}
}

\newcommand{\resumeSubHeadingListStart}{\begin{itemize}[leftmargin=0.15in,label={}]}
\newcommand{\resumeSubHeadingListEnd}{\end{itemize}}
\newcommand{\resumeItemListStart}{\begin{itemize}[leftmargin=0.18in,itemsep=1pt,topsep=2pt]}
\newcommand{\resumeItemListEnd}{\end{itemize}\vspace{-2pt}}

% ============================================================
\begin{document}
% ============================================================

%----------HEADING----------

\begin{center}

  {\LARGE \scshape \color{primary} Stéphane KPOVIESSI} \\ \vspace{3pt}

  {\large \textbf{Data \& AI Engineer}} \\ \vspace{4pt}

  \small
  Île-de-France (Cergy)
  $|$
  +33\,7\,48\,59\,47\,13
  $|$
  \href{mailto:oastephanekpoviessi@gmail.com}{oastephanekpoviessi@gmail.com}

  \vspace{2pt}

  \href{https://linkedin.com/in/stephanekpoviessi}{linkedin.com/in/stephanekpoviessi}
  $|$
  \href{https://github.com/EngineerProjects}{github.com/EngineerProjects}
  $|$
  \href{https://kpoviessi-stephane.vercel.app}{kpoviessi-stephane.vercel.app}

\end{center}

%-----------PROFIL-----------

\section{Profil}

\small{
  Ingénieur Data \& IA avec une expérience en conception de pipelines ETL industriels,
  architecture cloud moderne et développement de systèmes IA (RAG, agents LLM, orchestration).
  À l'aise aussi bien sur les problématiques data engineering classiques que sur les
  sujets IA appliqués en production. Recherche une équipe technique ambitieuse, idéalement
  dans une startup ou scale-up orientée data ou IA.
}

%-----------EXPERIENCES-----------

\section{Expériences Professionnelles}

\resumeSubHeadingListStart

  \resumeSubheading
    {Data Engineer}{Nov. 2025 -- Avr. 2026}
    {Allianz France}{Paris}

  \resumeItemListStart
    \resumeItem{
      Industrialisation de pipelines ETL PySpark sur Azure pour la migration d'un DataMart
      Assurance depuis un environnement legacy SAS, dans le cadre d'une modernisation à
      grande échelle. % @data_engineer
    }
    \resumeItem{
      Mise en place d'une architecture Médaillon (Bronze/Silver/Gold) avec plus de 150 règles
      métier implémentées et préparation à l'intégration Databricks. % @data_engineer
    }
    \resumeItem{
      Développement de pipelines de contrôle qualité et traçabilité des données clients
      avec validation automatique, tests de parité SAS/Python et interface de suivi. % @data_engineer @cost
    }
  \resumeItemListEnd

  \resumeSubheading
    {Machine Learning Engineer
      {\small
        \href{https://github.com/EngineerProjects/sylva3D}{[Sylva3D]}
        \href{https://github.com/EngineerProjects/Sylva3dGUI}{[Sylva3DGUI]}
      }
    }
    {Mai 2024 -- Août 2024}
    {Groupe Sylvagreg}{Lille}

  \resumeItemListStart
    \resumeItem{
      Développement d'un pipeline IA de génération 3D combinant diffusion models,
      reconstruction multi-vues et Computer Vision (PyTorch). % @ml_engineer @genai
    }
    \resumeItem{
      Déploiement de workflows Computer Vision sur AWS et développement d'une
      application desktop Python pour la visualisation des résultats. % @ml_engineer @backend
    }
  \resumeItemListEnd

\resumeSubHeadingListEnd

%-----------PROJECTS-----------

\section{Projets Techniques}

\resumeSubHeadingListStart

  \resumeProjectHeading
    {\textbf{BI Retail Solution} $|$
      \emph{Python, Airflow, PostgreSQL, Power BI, Docker}}
    {\href{https://github.com/EngineerProjects/BI_Retail}{GitHub}}

  \resumeItemListStart
    \resumeItem{
      Pipeline ETL retail complet avec orchestration Airflow, stockage PostgreSQL,
      dashboards Power BI et automatisation du reporting métier. % @data_engineer
    }
  \resumeItemListEnd

  \resumeProjectHeading
    {\textbf{Tech Watch Agent} $|$
      \emph{LangGraph, FastAPI, PostgreSQL, pgvector, Docker}}
    {\href{https://github.com/EngineerProjects/tech-watch-agent}{GitHub}}

  \resumeItemListStart
    \resumeItem{
      Plateforme multi-agents de veille technologique : workflows LLM, mémoire vectorielle
      (pgvector), génération automatisée de rapports et API REST. % @genai @rag
    }
  \resumeItemListEnd

  \resumeProjectHeading
    {\textbf{Nexus Engine} $|$
      \emph{Go, gRPC, SSE, SQLite, LLM Orchestration}}
    {} % projet privé, pas de lien GitHub

  \resumeItemListStart
    \resumeItem{
      Runtime IA orienté agents en Go : orchestration de workflows LLM, exécution de tools,
      streaming temps réel SSE/gRPC, gestion persistante de sessions, multi-providers. % @genai @backend
    }
    \resumeItem{
      Architecture modulaire conçue pour des systèmes IA interactifs, extensibles et
      orientés tooling -- production-ready from scratch. % @genai @backend
    }
  \resumeItemListEnd

\resumeSubHeadingListEnd

%-----------SKILLS-----------

\section{Compétences Techniques}

\begin{itemize}[leftmargin=0.15in,label={}]
  \small{\item{
    \textbf{Langages :} Python, SQL, Go \\
    \textbf{Data Engineering :} PySpark, Airflow, ETL/ELT, PostgreSQL, Architecture Médaillon, dbt \\
    \textbf{IA appliquée :} LangChain, LangGraph, RAG, systèmes multi-agents, orchestration LLM \\
    \textbf{Cloud \& Infra :} Azure, Azure Databricks, AWS, Docker, Git, Linux \\
    \textbf{Visualisation :} Power BI, Streamlit
    \vspace{4pt}

    {\color{primary}\textbf{Langues :}}
    Français, Anglais (B2), Yoruba, Fon, Goun (natifs)
  }}
\end{itemize}

%-----------FORMATION-----------

\section{Formation}

\resumeSubHeadingListStart

  \resumeSubheading
    {JUNIA ISEN}{Lille}
    {Diplôme d'Ingénieur --- Big Data \& Data Science}{2026}

  \resumeSubheading
    {JUNIA ISEN}{Rabat}
    {Licence --- Électronique}{2023}

\resumeSubHeadingListEnd

% ============================================================
\end{document}
% ============================================================
